% =====================================================================
%  labfisica.tex — macros propias del proyecto
%  Se carga desde el preámbulo de main.tex con % =====================================================================
%  labfisica.tex — macros propias del proyecto
%  Se carga desde el preámbulo de main.tex con % =====================================================================
%  labfisica.tex — macros propias del proyecto
%  Se carga desde el preámbulo de main.tex con % =====================================================================
%  labfisica.tex — macros propias del proyecto
%  Se carga desde el preámbulo de main.tex con \input{styles/labfisica}
% =====================================================================

\usepackage{keyval}

% --- Metadatos de una unidad ---------------------------------------
% Una "unidad" es una sesión (chapters/contents/sessio_XX/) o un
% cuadernillo (quaderns/XXX/). Las dos llevan su meta.tex:
%
%   \sessioMeta{
%     numero    = 1,
%     titol     = {Títol en català},
%     data      = {2026-09-14},
%     trimestre = {1},
%     blocs     = {materia},
%     estat     = {publicat},
%     solucions = {ocult},
%     resum     = {Una línia per al índex de la web.}
%   }
%
% El mismo fichero lo leerá el generador web, así que la fecha va
% siempre en formato ISO (aaaa-mm-dd) y se formatea al imprimir.
%
%   estat      publicat | esborrany   -> si se renderiza y se publica
%   solucions  ocult    | publicat    -> si el índice enlaza el PDF de solucions
%   trimestre  1 | 2 | 3              -> para agrupar en cuadernillos
%   resum                             -> una línea, para el índice de la web
%   numero                            -> vacío en un cuadernillo: no es una sessió
%   pdf                               -> ver abajo
%
% CAMPO pdf: para unidades SIN fuente LaTeX, maquetadas a mano (Affinity).
%
%   pdf = {quimica_t1.pdf}
%
% Si está presente, no hay nada que compilar: ese PDF es el original y es
% irrecuperable. El generador web lo publica tal cual y el .gitignore tiene
% una excepción para versionarlo. Ver tex/quaderns/README.md.

\makeatletter

\define@key{sessio}{numero}{\def\sessioNumero{#1}}
\define@key{sessio}{titol} {\def\sessioTitol{#1}}
\define@key{sessio}{data}  {\def\sessioData{#1}}
\define@key{sessio}{blocs} {\def\sessioBlocs{#1}}
\define@key{sessio}{estat} {\def\sessioEstat{#1}}
\define@key{sessio}{trimestre}{\def\sessioTrimestre{#1}}
\define@key{sessio}{solucions}{\def\sessioSolucions{#1}}
\define@key{sessio}{resum}    {\def\sessioResum{#1}}
\define@key{sessio}{pdf}      {\def\sessioPdf{#1}}

\newcommand{\sessioMeta}[1]{\setkeys{sessio}{#1}}

% 2026-09-14 --> 14/09/2026
\def\lf@datailegible#1-#2-#3\@nil{#3/#2/#1}
\newcommand{\lf@data}{\expandafter\lf@datailegible\sessioData\@nil}

\newcommand{\lf@publicat}{publicat}
\newcommand{\lf@buit}{}

% Valores por defecto. Se aplican antes de leer cada meta.tex, para que
% una clave ausente no herede el valor de la unidad anterior.
\newcommand{\lf@reset}{%
  \def\sessioNumero{}\def\sessioTitol{}\def\sessioData{0000-00-00}%
  \def\sessioBlocs{}\def\sessioEstat{esborrany}%
  \def\sessioTrimestre{}\def\sessioSolucions{ocult}\def\sessioResum{}%
  \def\sessioPdf{}%
}

% --- Carga de ficheros ---------------------------------------------
% Reciben la ruta COMPLETA sin extensión, relativa al directorio de
% compilación (tex/). Así valen igual para sesiones y cuadernillos.
\newcommand{\lf@obl}[1]{%
  \IfFileExists{#1.tex}%
    {\input{#1.tex}}%
    {\GenericWarning{}{labfisica: falta el fichero obligatorio #1.tex}}%
}
\newcommand{\lf@opc}[1]{%
  \IfFileExists{#1.tex}{\input{#1.tex}}{}%
}

% Las dos formas de citar una imagen valen: foto.jpg (histórico) y
% images/foto.jpg.
\newcommand{\lf@imatges}[1]{\graphicspath{{#1/images/}{#1/}{}}}

% Cabecera: "Sessió 1. Títol" si hay número; solo "Títol" si no lo hay
% (un cuadernillo no es una sessió).
\newcommand{\lf@capcalera}{%
  \ifx\sessioNumero\lf@buit
    \chapter{\sessioTitol}%
  \else
    \chapter{Sessió \sessioNumero. \sessioTitol}%
  \fi
}

% --- \sessio{sessio_01} --------------------------------------------
% Monta una sesión entera dentro del dossier: cabecera y partes en orden.
% Las sesiones con estat distinto de "publicat" NO se renderizan.
\newcommand{\sessio}[1]{%
  \lf@reset
  \lf@obl{chapters/contents/#1/meta}%
  \ifx\sessioEstat\lf@publicat
    \lf@imatges{chapters/contents/#1}%
    \lf@capcalera
    \lf@obl{chapters/contents/#1/continguts}%
    \lf@obl{chapters/contents/#1/activitats}%
    \lf@opc{chapters/contents/#1/investiga}%
  \else
    \GenericWarning{}{labfisica: sesion #1 omitida (estat=\sessioEstat)}%
  \fi
}

% --- \lfUnitat{ruta}{parte} ----------------------------------------
% Compila UNA parte suelta, con su cabecera y sus imágenes resueltas.
% La ruta es relativa a tex/, así que sirve para los dos sitios:
%
%   \lfUnitat{chapters/contents/sessio_01}{solucions}
%   \lfUnitat{quaderns/vectors}{continguts}
%
% A diferencia de \sessio, NO mira el estat: es la herramienta para
% revisar un borrador en local antes de publicarlo.
\newcommand{\lfUnitat}[2]{%
  \lf@reset
  \lf@obl{#1/meta}%
  \lf@imatges{#1}%
  \lf@capcalera
  \lf@obl{#1/#2}%
}

\makeatother

% =====================================================================

\usepackage{keyval}

% --- Metadatos de una unidad ---------------------------------------
% Una "unidad" es una sesión (chapters/contents/sessio_XX/) o un
% cuadernillo (quaderns/XXX/). Las dos llevan su meta.tex:
%
%   \sessioMeta{
%     numero    = 1,
%     titol     = {Títol en català},
%     data      = {2026-09-14},
%     trimestre = {1},
%     blocs     = {materia},
%     estat     = {publicat},
%     solucions = {ocult},
%     resum     = {Una línia per al índex de la web.}
%   }
%
% El mismo fichero lo leerá el generador web, así que la fecha va
% siempre en formato ISO (aaaa-mm-dd) y se formatea al imprimir.
%
%   estat      publicat | esborrany   -> si se renderiza y se publica
%   solucions  ocult    | publicat    -> si el índice enlaza el PDF de solucions
%   trimestre  1 | 2 | 3              -> para agrupar en cuadernillos
%   resum                             -> una línea, para el índice de la web
%   numero                            -> vacío en un cuadernillo: no es una sessió
%   pdf                               -> ver abajo
%
% CAMPO pdf: para unidades SIN fuente LaTeX, maquetadas a mano (Affinity).
%
%   pdf = {quimica_t1.pdf}
%
% Si está presente, no hay nada que compilar: ese PDF es el original y es
% irrecuperable. El generador web lo publica tal cual y el .gitignore tiene
% una excepción para versionarlo. Ver tex/quaderns/README.md.

\makeatletter

\define@key{sessio}{numero}{\def\sessioNumero{#1}}
\define@key{sessio}{titol} {\def\sessioTitol{#1}}
\define@key{sessio}{data}  {\def\sessioData{#1}}
\define@key{sessio}{blocs} {\def\sessioBlocs{#1}}
\define@key{sessio}{estat} {\def\sessioEstat{#1}}
\define@key{sessio}{trimestre}{\def\sessioTrimestre{#1}}
\define@key{sessio}{solucions}{\def\sessioSolucions{#1}}
\define@key{sessio}{resum}    {\def\sessioResum{#1}}
\define@key{sessio}{pdf}      {\def\sessioPdf{#1}}

\newcommand{\sessioMeta}[1]{\setkeys{sessio}{#1}}

% 2026-09-14 --> 14/09/2026
\def\lf@datailegible#1-#2-#3\@nil{#3/#2/#1}
\newcommand{\lf@data}{\expandafter\lf@datailegible\sessioData\@nil}

\newcommand{\lf@publicat}{publicat}
\newcommand{\lf@buit}{}

% Valores por defecto. Se aplican antes de leer cada meta.tex, para que
% una clave ausente no herede el valor de la unidad anterior.
\newcommand{\lf@reset}{%
  \def\sessioNumero{}\def\sessioTitol{}\def\sessioData{0000-00-00}%
  \def\sessioBlocs{}\def\sessioEstat{esborrany}%
  \def\sessioTrimestre{}\def\sessioSolucions{ocult}\def\sessioResum{}%
  \def\sessioPdf{}%
}

% --- Carga de ficheros ---------------------------------------------
% Reciben la ruta COMPLETA sin extensión, relativa al directorio de
% compilación (tex/). Así valen igual para sesiones y cuadernillos.
\newcommand{\lf@obl}[1]{%
  \IfFileExists{#1.tex}%
    {\input{#1.tex}}%
    {\GenericWarning{}{labfisica: falta el fichero obligatorio #1.tex}}%
}
\newcommand{\lf@opc}[1]{%
  \IfFileExists{#1.tex}{\input{#1.tex}}{}%
}

% Las dos formas de citar una imagen valen: foto.jpg (histórico) y
% images/foto.jpg.
\newcommand{\lf@imatges}[1]{\graphicspath{{#1/images/}{#1/}{}}}

% Cabecera: "Sessió 1. Títol" si hay número; solo "Títol" si no lo hay
% (un cuadernillo no es una sessió).
\newcommand{\lf@capcalera}{%
  \ifx\sessioNumero\lf@buit
    \chapter{\sessioTitol}%
  \else
    \chapter{Sessió \sessioNumero. \sessioTitol}%
  \fi
}

% --- \sessio{sessio_01} --------------------------------------------
% Monta una sesión entera dentro del dossier: cabecera y partes en orden.
% Las sesiones con estat distinto de "publicat" NO se renderizan.
\newcommand{\sessio}[1]{%
  \lf@reset
  \lf@obl{chapters/contents/#1/meta}%
  \ifx\sessioEstat\lf@publicat
    \lf@imatges{chapters/contents/#1}%
    \lf@capcalera
    \lf@obl{chapters/contents/#1/continguts}%
    \lf@obl{chapters/contents/#1/activitats}%
    \lf@opc{chapters/contents/#1/investiga}%
  \else
    \GenericWarning{}{labfisica: sesion #1 omitida (estat=\sessioEstat)}%
  \fi
}

% --- \lfUnitat{ruta}{parte} ----------------------------------------
% Compila UNA parte suelta, con su cabecera y sus imágenes resueltas.
% La ruta es relativa a tex/, así que sirve para los dos sitios:
%
%   \lfUnitat{chapters/contents/sessio_01}{solucions}
%   \lfUnitat{quaderns/vectors}{continguts}
%
% A diferencia de \sessio, NO mira el estat: es la herramienta para
% revisar un borrador en local antes de publicarlo.
\newcommand{\lfUnitat}[2]{%
  \lf@reset
  \lf@obl{#1/meta}%
  \lf@imatges{#1}%
  \lf@capcalera
  \lf@obl{#1/#2}%
}

\makeatother

% =====================================================================

\usepackage{keyval}

% --- Metadatos de una unidad ---------------------------------------
% Una "unidad" es una sesión (chapters/contents/sessio_XX/) o un
% cuadernillo (quaderns/XXX/). Las dos llevan su meta.tex:
%
%   \sessioMeta{
%     numero    = 1,
%     titol     = {Títol en català},
%     data      = {2026-09-14},
%     trimestre = {1},
%     blocs     = {materia},
%     estat     = {publicat},
%     solucions = {ocult},
%     resum     = {Una línia per al índex de la web.}
%   }
%
% El mismo fichero lo leerá el generador web, así que la fecha va
% siempre en formato ISO (aaaa-mm-dd) y se formatea al imprimir.
%
%   estat      publicat | esborrany   -> si se renderiza y se publica
%   solucions  ocult    | publicat    -> si el índice enlaza el PDF de solucions
%   trimestre  1 | 2 | 3              -> para agrupar en cuadernillos
%   resum                             -> una línea, para el índice de la web
%   numero                            -> vacío en un cuadernillo: no es una sessió
%   pdf                               -> ver abajo
%
% CAMPO pdf: para unidades SIN fuente LaTeX, maquetadas a mano (Affinity).
%
%   pdf = {quimica_t1.pdf}
%
% Si está presente, no hay nada que compilar: ese PDF es el original y es
% irrecuperable. El generador web lo publica tal cual y el .gitignore tiene
% una excepción para versionarlo. Ver tex/quaderns/README.md.

\makeatletter

\define@key{sessio}{numero}{\def\sessioNumero{#1}}
\define@key{sessio}{titol} {\def\sessioTitol{#1}}
\define@key{sessio}{data}  {\def\sessioData{#1}}
\define@key{sessio}{blocs} {\def\sessioBlocs{#1}}
\define@key{sessio}{estat} {\def\sessioEstat{#1}}
\define@key{sessio}{trimestre}{\def\sessioTrimestre{#1}}
\define@key{sessio}{solucions}{\def\sessioSolucions{#1}}
\define@key{sessio}{resum}    {\def\sessioResum{#1}}
\define@key{sessio}{pdf}      {\def\sessioPdf{#1}}

\newcommand{\sessioMeta}[1]{\setkeys{sessio}{#1}}

% 2026-09-14 --> 14/09/2026
\def\lf@datailegible#1-#2-#3\@nil{#3/#2/#1}
\newcommand{\lf@data}{\expandafter\lf@datailegible\sessioData\@nil}

\newcommand{\lf@publicat}{publicat}
\newcommand{\lf@buit}{}

% Valores por defecto. Se aplican antes de leer cada meta.tex, para que
% una clave ausente no herede el valor de la unidad anterior.
\newcommand{\lf@reset}{%
  \def\sessioNumero{}\def\sessioTitol{}\def\sessioData{0000-00-00}%
  \def\sessioBlocs{}\def\sessioEstat{esborrany}%
  \def\sessioTrimestre{}\def\sessioSolucions{ocult}\def\sessioResum{}%
  \def\sessioPdf{}%
}

% --- Carga de ficheros ---------------------------------------------
% Reciben la ruta COMPLETA sin extensión, relativa al directorio de
% compilación (tex/). Así valen igual para sesiones y cuadernillos.
\newcommand{\lf@obl}[1]{%
  \IfFileExists{#1.tex}%
    {\input{#1.tex}}%
    {\GenericWarning{}{labfisica: falta el fichero obligatorio #1.tex}}%
}
\newcommand{\lf@opc}[1]{%
  \IfFileExists{#1.tex}{\input{#1.tex}}{}%
}

% Las dos formas de citar una imagen valen: foto.jpg (histórico) y
% images/foto.jpg.
\newcommand{\lf@imatges}[1]{\graphicspath{{#1/images/}{#1/}{}}}

% Cabecera: "Sessió 1. Títol" si hay número; solo "Títol" si no lo hay
% (un cuadernillo no es una sessió).
\newcommand{\lf@capcalera}{%
  \ifx\sessioNumero\lf@buit
    \chapter{\sessioTitol}%
  \else
    \chapter{Sessió \sessioNumero. \sessioTitol}%
  \fi
}

% --- \sessio{sessio_01} --------------------------------------------
% Monta una sesión entera dentro del dossier: cabecera y partes en orden.
% Las sesiones con estat distinto de "publicat" NO se renderizan.
\newcommand{\sessio}[1]{%
  \lf@reset
  \lf@obl{chapters/contents/#1/meta}%
  \ifx\sessioEstat\lf@publicat
    \lf@imatges{chapters/contents/#1}%
    \lf@capcalera
    \lf@obl{chapters/contents/#1/continguts}%
    \lf@obl{chapters/contents/#1/activitats}%
    \lf@opc{chapters/contents/#1/investiga}%
  \else
    \GenericWarning{}{labfisica: sesion #1 omitida (estat=\sessioEstat)}%
  \fi
}

% --- \lfUnitat{ruta}{parte} ----------------------------------------
% Compila UNA parte suelta, con su cabecera y sus imágenes resueltas.
% La ruta es relativa a tex/, así que sirve para los dos sitios:
%
%   \lfUnitat{chapters/contents/sessio_01}{solucions}
%   \lfUnitat{quaderns/vectors}{continguts}
%
% A diferencia de \sessio, NO mira el estat: es la herramienta para
% revisar un borrador en local antes de publicarlo.
\newcommand{\lfUnitat}[2]{%
  \lf@reset
  \lf@obl{#1/meta}%
  \lf@imatges{#1}%
  \lf@capcalera
  \lf@obl{#1/#2}%
}

\makeatother

% =====================================================================

\usepackage{keyval}

% --- Metadatos de una unidad ---------------------------------------
% Una "unidad" es una sesión (chapters/contents/sessio_XX/) o un
% cuadernillo (quaderns/XXX/). Las dos llevan su meta.tex:
%
%   \sessioMeta{
%     numero    = 1,
%     titol     = {Títol en català},
%     data      = {2026-09-14},
%     trimestre = {1},
%     blocs     = {materia},
%     estat     = {publicat},
%     solucions = {ocult},
%     resum     = {Una línia per al índex de la web.}
%   }
%
% El mismo fichero lo leerá el generador web, así que la fecha va
% siempre en formato ISO (aaaa-mm-dd) y se formatea al imprimir.
%
%   estat      publicat | esborrany   -> si se renderiza y se publica
%   solucions  ocult    | publicat    -> si el índice enlaza el PDF de solucions
%   trimestre  1 | 2 | 3              -> para agrupar en cuadernillos
%   resum                             -> una línea, para el índice de la web
%   numero                            -> vacío en un cuadernillo: no es una sessió
%   pdf                               -> ver abajo
%
% CAMPO pdf: para unidades SIN fuente LaTeX, maquetadas a mano (Affinity).
%
%   pdf = {quimica_t1.pdf}
%
% Si está presente, no hay nada que compilar: ese PDF es el original y es
% irrecuperable. El generador web lo publica tal cual y el .gitignore tiene
% una excepción para versionarlo. Ver tex/quaderns/README.md.

\makeatletter

\define@key{sessio}{numero}{\def\sessioNumero{#1}}
\define@key{sessio}{titol} {\def\sessioTitol{#1}}
\define@key{sessio}{data}  {\def\sessioData{#1}}
\define@key{sessio}{blocs} {\def\sessioBlocs{#1}}
\define@key{sessio}{estat} {\def\sessioEstat{#1}}
\define@key{sessio}{trimestre}{\def\sessioTrimestre{#1}}
\define@key{sessio}{solucions}{\def\sessioSolucions{#1}}
\define@key{sessio}{resum}    {\def\sessioResum{#1}}
\define@key{sessio}{pdf}      {\def\sessioPdf{#1}}

\newcommand{\sessioMeta}[1]{\setkeys{sessio}{#1}}

% 2026-09-14 --> 14/09/2026
\def\lf@datailegible#1-#2-#3\@nil{#3/#2/#1}
\newcommand{\lf@data}{\expandafter\lf@datailegible\sessioData\@nil}

\newcommand{\lf@publicat}{publicat}
\newcommand{\lf@buit}{}

% Valores por defecto. Se aplican antes de leer cada meta.tex, para que
% una clave ausente no herede el valor de la unidad anterior.
\newcommand{\lf@reset}{%
  \def\sessioNumero{}\def\sessioTitol{}\def\sessioData{0000-00-00}%
  \def\sessioBlocs{}\def\sessioEstat{esborrany}%
  \def\sessioTrimestre{}\def\sessioSolucions{ocult}\def\sessioResum{}%
  \def\sessioPdf{}%
}

% --- Carga de ficheros ---------------------------------------------
% Reciben la ruta COMPLETA sin extensión, relativa al directorio de
% compilación (tex/). Así valen igual para sesiones y cuadernillos.
\newcommand{\lf@obl}[1]{%
  \IfFileExists{#1.tex}%
    {\input{#1.tex}}%
    {\GenericWarning{}{labfisica: falta el fichero obligatorio #1.tex}}%
}
\newcommand{\lf@opc}[1]{%
  \IfFileExists{#1.tex}{\input{#1.tex}}{}%
}

% Las dos formas de citar una imagen valen: foto.jpg (histórico) y
% images/foto.jpg.
\newcommand{\lf@imatges}[1]{\graphicspath{{#1/images/}{#1/}{}}}

% Cabecera: "Sessió 1. Títol" si hay número; solo "Títol" si no lo hay
% (un cuadernillo no es una sessió).
\newcommand{\lf@capcalera}{%
  \ifx\sessioNumero\lf@buit
    \chapter{\sessioTitol}%
  \else
    \chapter{Sessió \sessioNumero. \sessioTitol}%
  \fi
}

% --- \sessio{sessio_01} --------------------------------------------
% Monta una sesión entera dentro del dossier: cabecera y partes en orden.
% Las sesiones con estat distinto de "publicat" NO se renderizan.
\newcommand{\sessio}[1]{%
  \lf@reset
  \lf@obl{chapters/contents/#1/meta}%
  \ifx\sessioEstat\lf@publicat
    \lf@imatges{chapters/contents/#1}%
    \lf@capcalera
    \lf@obl{chapters/contents/#1/continguts}%
    \lf@obl{chapters/contents/#1/activitats}%
    \lf@opc{chapters/contents/#1/investiga}%
  \else
    \GenericWarning{}{labfisica: sesion #1 omitida (estat=\sessioEstat)}%
  \fi
}

% --- \lfUnitat{ruta}{parte} ----------------------------------------
% Compila UNA parte suelta, con su cabecera y sus imágenes resueltas.
% La ruta es relativa a tex/, así que sirve para los dos sitios:
%
%   \lfUnitat{chapters/contents/sessio_01}{solucions}
%   \lfUnitat{quaderns/vectors}{continguts}
%
% A diferencia de \sessio, NO mira el estat: es la herramienta para
% revisar un borrador en local antes de publicarlo.
\newcommand{\lfUnitat}[2]{%
  \lf@reset
  \lf@obl{#1/meta}%
  \lf@imatges{#1}%
  \lf@capcalera
  \lf@obl{#1/#2}%
}

\makeatother
